% Work/MAT | 1. rocnik | 34: Slozene vyrazy s mocninami a odmocninami
% Navaznost: 25--30 mocniny, 31 castecne odmocnovani,
% 32 usmernovani jednoduchych zlomku, 33 sdruzene vyrazy.
% Zdroje: ucebnice Mocniny a Odmocniny.tex (Caste chyby, Odmocniny,
% Usmernovani zlomku); Realisticky.cz 020712--020714 Pocitani s odmocninami.
% Zcela nove formulovane priklady. Dokument je priprava pro ucitele,
% 45minutove jadro + vyslovne oznacene volitelne cviceni.
\documentclass[11pt,a4paper]{article}
% Spolecny styl pro ucitelske pripravy MAT (Work/MAT)
% Verze 1.1, 2026-10-23
\NeedsTeXFormat{LaTeX2e}
\ProvidesPackage{priprava}[2026/10/23 v1.1 Sablona ucitelskych priprav MAT]

% Zakladni typografie a cestina
\RequirePackage[T1]{fontenc}
\RequirePackage[utf8]{inputenc}
\RequirePackage[czech]{babel}
\RequirePackage{lmodern}
\RequirePackage{microtype}
\RequirePackage{amsmath,amssymb}
\RequirePackage{xcolor}
\RequirePackage{enumitem}
\RequirePackage[most]{tcolorbox}
\RequirePackage{tabularx,booktabs,array}
\RequirePackage{fancyhdr}
\RequirePackage{needspace}
\RequirePackage{geometry}
\RequirePackage[hidelinks]{hyperref}

\geometry{a4paper,top=17mm,bottom=18mm,left=23mm,right=23mm}
\definecolor{notegray}{gray}{0.39}
\setlength{\parindent}{0pt}
\setlength{\parskip}{2.5pt plus 1pt minus 0.5pt}
\setlist[itemize]{topsep=2pt,itemsep=1pt,parsep=0pt,leftmargin=1.45em}
\setlist[enumerate]{topsep=2pt,itemsep=2pt,parsep=0pt,leftmargin=1.65em}
\setcounter{secnumdepth}{0}
\renewcommand{\arraystretch}{1.12}

% Zahlavi: vlevo rocnik | globalni tema, vpravo cislo: nazev hodiny
% Zapati zustava zcela prazdne.
\pagestyle{fancy}
\fancyhf{}
\renewcommand{\headrulewidth}{0.35pt}
\renewcommand{\footrulewidth}{0pt}
\setlength{\headheight}{14pt}
\newcommand{\nastavPripravu}[4]{%
  \fancyhead[L]{\small #1\enspace\textbar\enspace #2}%
  \fancyhead[R]{\small #3: #4}%
}

% Minimalisticke cernobile boxy. Kazda definice zvlast.
\tcbset{commonlesson/.style={enhanced,breakable=false,colback=white,colframe=black,
  colbacktitle=white,coltitle=black,boxrule=0.45pt,
  arc=0.8mm,outer arc=0.8mm,left=3mm,right=3mm,top=1.6mm,bottom=1.6mm,
  toptitle=1mm,bottomtitle=0.7mm,fonttitle=\bfseries,
  attach boxed title to top left={xshift=2.5mm,yshift=-1.15mm},
  boxed title style={colback=white,colframe=white,boxrule=0pt,left=1mm,right=1mm}}}
\newtcolorbox{zadaniBox}[1][]{commonlesson,title={Zadání},#1}
\newtcolorbox{prikladBox}[1][]{commonlesson,title={Příklad},#1}
\newtcolorbox{definiceBox}[1][]{commonlesson,title={Definice},#1}

% Spolecny box pro zadani a jeho reseni, bez mezery mezi castmi.
% Pouziti:
%   \begin{prikladReseniBox}
%     Zadani ulohy.
%     \tcblower
%     \begin{align*} ... \end{align*}
%   \end{prikladReseniBox}
% Spodni cast ma jemny sedy podklad a delici linku; lze v ni sazet
% vice radku matematiky. Obe casti maji jen jeden spolecny ramecek.
\tcbset{solutionlesson/.style={
  commonlesson,
  enhanced jigsaw,
  breakable,
  colbacklower=black!3,
  segmentation style={solid,draw=black!25,line width=0.35pt},
  fontlower=\small\color{notegray},
  before lower={\noindent\textbf{Řešení:}\enspace},
  leftlower=3mm,rightlower=3mm
}}
\newtcolorbox{prikladReseniBox}[1][]{solutionlesson,title={Příklad},#1}
\newtcolorbox{zadaniReseniBox}[1][]{solutionlesson,title={Zadání},#1}

% Pomucky pro autora pripravy: sedive poznamky; \reseni pro starsi pripravy
\newcommand{\poznamka}[1]{\par\begingroup\small\color{notegray}\itshape
  \textnormal{\textbf{Poznámka pro učitele.}}\enspace #1\par\endgroup}
\newcommand{\reseni}[1]{\par\begingroup\small\color{notegray}
  \textbf{Očekávané řešení:}\enspace #1\par\endgroup}
\newcommand{\blok}[1]{\needspace{5\baselineskip}\subsection*{#1}}
\newcommand{\vyrok}[1]{\uv{#1}}

\usepackage{multicol}
\nastavPripravu{1. ročník}{Mocniny a odmocniny}{34}{Složené výrazy}
\begin{document}

\poznamka{\textbf{45minutové jádro:} 4 min návrat k pravidlům;
8 min Příklad 1 a strategie; 8 min Samostatná práce 1 a)--d);
8 min Příklad 2 a Samostatná práce 2 a)--c);
9 min Příklad 3 a Samostatná práce 3 a)--c);
5 min společný rozbor chyb; 3 min výstupní otázka a domácí úkol.
\textbf{Nadstavba:} další položky, Příklad 4, rozšiřující příklady.
Nejde o závod o počet příkladů, ale o správně zvolený postup.}

\blok{KISS v praxi: nejprve plán, potom výpočet}
V minulých hodinách jsme zvládli pravidla pro mocniny, částečné
odmocňování, usměrňování jmenovatele i sdružené výrazy.
Dnes \textbf{nepřidáváme nový vzorec}: učíme se rozhodnout, co použít
\emph{nejdříve}, aby byl výpočet co nejkratší.

\begin{definiceBox}[title={Doporučená strategie pro složený výraz}]
\begin{enumerate}[label=\arabic*.]
\item \textbf{Podmínky:} označ jmenovatele a sudé odmocniny; žádný jmenovatel nesmí být nula.
\item \textbf{Zjednoduš:} částečně odmocni, sjednoť stejné základy a sluč podobné odmocniny.
\item \textbf{Až potom} násob, umocňuj dvojčlen či případně usměrňuj.
\item \textbf{Zkontroluj:} znaménka, podmínky a zachování hodnoty výrazu.
\end{enumerate}
\end{definiceBox}
\poznamka{Body strategie nejsou čtyři povinné početní operace; pokud
po prvním kroku zlomky zmizí, dále neusměrňujeme. Zeptej se:
\uv{Kterou úpravu bys v tomto příkladu provedl jako první a proč?}}

\blok{1. Když odmocniny nejprve zjednodušíme, vyjde vše samo}
\begin{prikladReseniBox}[breakable=false,title={Příklad 1: co půjde zkrátit bez násobení?}]
Uprav $\displaystyle\frac{\sqrt{72}-\sqrt{18}}{\sqrt2}$ a
$\displaystyle\frac{\sqrt{48}+\sqrt{27}}{\sqrt3}$.
\tcblower
\[
\frac{6\sqrt2-3\sqrt2}{\sqrt2}=\boxed{3},\qquad
\frac{4\sqrt3+3\sqrt3}{\sqrt3}=\boxed{7}.
\]
\textbf{Strategie:} nejprve vytkneme vhodný čtverec pod odmocninou,
potom sečteme koeficienty. Jmenovatele není nutné usměrňovat.
\end{prikladReseniBox}

\begin{zadaniReseniBox}[breakable=false,title={Samostatná práce 1: odmocni, slučuj, zjednoduš}]
Uprav přesně, bez desetinného přiblížení:
\begin{multicols}{2}
\begin{enumerate}[label=\alph*)]
\item $\sqrt{50}+\sqrt8-\sqrt{18}$;
\item $\displaystyle\frac{\sqrt{48}-\sqrt{12}}{\sqrt3}$;
\item $\displaystyle\frac{\sqrt{45}+\sqrt{20}}{\sqrt5}$;
\item $\displaystyle\frac{\sqrt[3]{54}-\sqrt[3]{16}}{\sqrt[3]2}$;
\item $\displaystyle\frac{(\sqrt6)^2+\sqrt{81}}{\sqrt3}$;
\item $\displaystyle\frac{\sqrt{98}-2\sqrt8}{\sqrt2}$.
\end{enumerate}
\end{multicols}
\tcblower
\begin{multicols}{2}
\begin{enumerate}[label=\alph*)]
\item $(5+2-3)\sqrt2=\boxed{4\sqrt2}$;
\item $(4\sqrt3-2\sqrt3)/\sqrt3=\boxed2$;
\item $(3\sqrt5+2\sqrt5)/\sqrt5=\boxed5$;
\item $(3\sqrt[3]2-2\sqrt[3]2)/\sqrt[3]2=\boxed1$;
\item $15/\sqrt3=\boxed{5\sqrt3}$;
\item $(7\sqrt2-4\sqrt2)/\sqrt2=\boxed3$.
\end{enumerate}
\end{multicols}
\end{zadaniReseniBox}
\poznamka{První čtyři položky tvoří povinné jádro, poslední dvě
slouží k diferenciaci. U d) zdůrazni rozdíl mezi druhou a třetí odmocninou.}

\blok{2. Sdružený výraz použij až tehdy, když opravdu pomůže}
\begin{prikladReseniBox}[breakable=false,title={Příklad 2: rozdíl dvou zlomků}]
Uprav $\displaystyle\frac2{\sqrt3-1}-\frac2{\sqrt3+1}$.
\tcblower
Usměrníme každý jmenovatel až po rozpoznání sdružených výrazů:
\[
\frac{2(\sqrt3+1)}{3-1}-\frac{2(\sqrt3-1)}{3-1}
=(\sqrt3+1)-(\sqrt3-1)=\boxed2.
\]
\textbf{Kontrola:} žádný původní jmenovatel není nula.
\end{prikladReseniBox}

\begin{zadaniReseniBox}[breakable=false,title={Samostatná práce 2: vyber vhodné pravidlo}]
\begin{multicols}{2}
\begin{enumerate}[label=\alph*)]
\item $\displaystyle\frac1{\sqrt5-2}-\frac1{\sqrt5+2}$;
\item $\displaystyle\frac2{\sqrt3+1}$;
\item $(\sqrt7+\sqrt3)(\sqrt7-\sqrt3)$;
\item $(\sqrt3-1)^2$;
\item $\displaystyle\frac{\sqrt2}{\sqrt3-\sqrt2}$;
\item $\displaystyle\frac1{\sqrt3+\sqrt2}$.
\end{enumerate}
\end{multicols}
\tcblower
\begin{multicols}{2}
\begin{enumerate}[label=\alph*)]
\item $(\sqrt5+2)-(\sqrt5-2)=\boxed4$;
\item $2(\sqrt3-1)/(3-1)=\boxed{\sqrt3-1}$;
\item $7-3=\boxed4$;
\item $3-2\sqrt3+1=\boxed{4-2\sqrt3}$;
\item $\sqrt2(\sqrt3+\sqrt2)/(3-2)=\boxed{\sqrt6+2}$;
\item $(\sqrt3-\sqrt2)/(3-2)=\boxed{\sqrt3-\sqrt2}$.
\end{enumerate}
\end{multicols}
\end{zadaniReseniBox}

\blok{3. Mocniny, odmocniny a podmínky v jednom výrazu}
\begin{prikladReseniBox}[breakable=false,title={Příklad 3: správné pořadí úprav}]
Uprav $\displaystyle\frac{(\sqrt{12}+\sqrt{27})^2}{3}$
a při $a>0$ výraz $\displaystyle\frac{a^{3/2}a^{-1/2}}{\sqrt{a^2}}$.
\tcblower
\[
\frac{(2\sqrt3+3\sqrt3)^2}{3}
=\frac{(5\sqrt3)^2}{3}=\boxed{25}.
\]
Pro $a>0$ je $\sqrt{a^2}=a$, proto
\[
\frac{a^{3/2-1/2}}{\sqrt{a^2}}=\frac aa=\boxed1.
\]
Bez podmínky $a>0$ bychom museli pracovat s $|a|$.
\end{prikladReseniBox}

\begin{zadaniReseniBox}[breakable=false,title={Samostatná práce 3: smíšené úpravy}]
\begin{enumerate}[label=\alph*)]
\item $\displaystyle\frac{\sqrt{27}-\sqrt{12}}{\sqrt3}$;
\item $\displaystyle\frac{27^{2/3}}{\sqrt9}$;
\item $\displaystyle\frac{(\sqrt8+\sqrt{18})^2}{2}$;
\item $\displaystyle\frac{\sqrt{x^2}}x$, $x\ne0$;
\item $\displaystyle\frac{(\sqrt a\,b)^2}{ab}$, $a,b>0$.
\end{enumerate}
\tcblower
\begin{enumerate}[label=\alph*)]
\item $(3\sqrt3-2\sqrt3)/\sqrt3=\boxed1$;
\item $9/3=\boxed3$;
\item $(2\sqrt2+3\sqrt2)^2/2=50/2=\boxed{25}$;
\item $\displaystyle\frac{|x|}x=\begin{cases}1,&x>0,\\-1,&x<0;\end{cases}$
\quad(nelze obecně napsat $1$);
\item $(ab^2)/(ab)=\boxed b$, při daných $a,b>0$.
\end{enumerate}
\end{zadaniReseniBox}
\poznamka{V povinném jádru stačí Samostatná práce 3 a)--c).
Další dvě úlohy jsou pro rychlé a testují zachování podmínek.}

\blok{4. Rozšiřující úlohy a práce s chybou}
\begin{prikladReseniBox}[breakable=false,title={Příklad 4: kdy se vyplatí jeden velký krok?}]
Uprav $\displaystyle
\frac{\sqrt{75}-\sqrt{12}}{\sqrt3+1}$.
\tcblower
Nejprve částečně odmocníme: $\sqrt{75}-\sqrt{12}=3\sqrt3$.
Teprve potom usměrníme:
\[
\frac{3\sqrt3}{\sqrt3+1}\cdot\frac{\sqrt3-1}{\sqrt3-1}
=\boxed{\frac{9-3\sqrt3}{2}}.
\]
\end{prikladReseniBox}

\begin{zadaniReseniBox}[breakable=false,title={Rozšiřující příklad 1: porovnej strategie}]
\begin{enumerate}[label=\alph*)]
\item $\displaystyle\frac{\sqrt{32}+\sqrt{18}}{\sqrt2}
\cdot\frac1{\sqrt5-2}$;
\item $\displaystyle\frac{(\sqrt3+\sqrt2)^2-(\sqrt3-\sqrt2)^2}{\sqrt6}$;
\item Pro $a,b>0$ uprav $\displaystyle\frac{(a^{1/2}b^{-1})^2}{a^{-1}b^{-3}}$.
\end{enumerate}
\tcblower
\begin{enumerate}[label=\alph*)]
\item $7/(\sqrt5-2)=\boxed{7(\sqrt5+2)}$;
\item Po odečtení čtverců vznikne $4\sqrt6/\sqrt6=\boxed4$;
\item $a^{1+1}b^{-2+3}=\boxed{a^2b}$.
\end{enumerate}
\end{zadaniReseniBox}

\begin{zadaniReseniBox}[breakable=false,title={Rozšiřující příklad 2: co je špatně?}]
Žák tvrdí:
\begin{enumerate}[label=\alph*)]
\item $(\sqrt5+\sqrt3)^2=5+3$;
\item $\sqrt{x^2}=x$ pro každé $x\in\mathbb R$;
\item $\sqrt{a+b}=\sqrt a+\sqrt b$ pro $a,b\ge0$.
\end{enumerate}
Oprav všechna tvrzení a uveď vhodný protipříklad.
\tcblower
\begin{enumerate}[label=\alph*)]
\item Chybí člen $2\sqrt{15}$; správně $8+2\sqrt{15}$;
\item správně $|x|$, například pro $x=-2$ je $\sqrt4=2$;
\item obecně neplatí, například $a=9$, $b=16$ dává $5\ne7$.
\end{enumerate}
\end{zadaniReseniBox}

\blok{Domácí procvičování}
\begin{zadaniReseniBox}[breakable=false,title={Domácí úkol 1: směs probraných operací}]
\begin{multicols}{2}
\begin{enumerate}[label=\alph*)]
\item $\displaystyle\frac{\sqrt{32}+\sqrt{18}}{\sqrt2}$;
\item $\sqrt{75}-2\sqrt{12}+\sqrt{27}$;
\item $\displaystyle\frac1{\sqrt5-1}$;
\item $\displaystyle\frac{2\sqrt3}{\sqrt3+1}$;
\item $\displaystyle\frac{(\sqrt{20}-\sqrt5)^2}{5}$;
\item $\displaystyle\frac{\sqrt[3]{128}-\sqrt[3]{16}}{\sqrt[3]2}$;
\item $16^{3/4}-81^{1/2}$;
\item $\displaystyle\frac{\sqrt{x^2}}x$, $x<0$.
\end{enumerate}
\end{multicols}
\tcblower
\begin{multicols}{2}
\begin{enumerate}[label=\alph*)]
\item $\boxed7$;
\item $(5-4+3)\sqrt3=\boxed{4\sqrt3}$;
\item $\boxed{(\sqrt5+1)/4}$;
\item $\boxed{3-\sqrt3}$;
\item $(2\sqrt5-\sqrt5)^2/5=\boxed1$;
\item $(4\sqrt[3]2-2\sqrt[3]2)/\sqrt[3]2=\boxed2$;
\item $8-9=\boxed{-1}$;
\item $|x|/x=\boxed{-1}$ pro $x<0$.
\end{enumerate}
\end{multicols}
\end{zadaniReseniBox}

\poznamka{\textbf{Rychlá výstupní otázka:} Proč je často výhodnější
nejprve částečně odmocnit a teprve potom usměrnit? Žák má odpovědět,
že se mohou zkrátit činitelé nebo sloučit členy a usměrňování již nemusí
být potřeba. V hodině 35 navážeme praktickou stránkou mocnin desítky.}
\end{document}
